\section{Experimental Setup}
\label{sec:setup}

\paragraph{Cluster and observability.}
All data were collected on \texttt{observit-cluster1}, a real nine-node RKE2 cluster (Kubernetes
v1.32.7) running the Online Boutique demo. Two observability stacks coexist: Prometheus/Grafana
for mature metrics and an OpenTelemetry collector for traces and logs, correlated through OTel
semantic conventions. The service set is fixed at $N=6$ (\texttt{frontend}, \texttt{cart},
\texttt{recommendation}, \texttt{ad}, \texttt{product-catalog}, \texttt{load-generator}); services
invisible on any modality were excluded to avoid systemic NaNs.

\paragraph{Fault injection.}
Faults are injected with Chaos~Mesh over 15 scenarios: 4 benign drifts
(\texttt{drift\_config\_change}, \texttt{drift\_rolling\_deploy}, \texttt{drift\_scale\_up},
\texttt{drift\_traffic\_ramp}) and 11 anomalies, including \texttt{faulty\_deploy\_overlap} which
realises the combined regime $\thetadn$. Each episode follows the
\emph{record\,$\to$\,build\,$\to$\,assemble} protocol: raw dumps are immutable ground truth,
features are recomputed offline, and datasets are split last.

\paragraph{Datasets.}
Table~\ref{tab:datasets} summarises the four datasets. We highlight one methodological pitfall:
the original temporal split of \texttt{ewat\_v4} left four scenarios entirely absent from the
training set, collapsing the independent-target AUROC to a trivial $0.5$. We therefore assembled
\texttt{ewat\_v4\_strat} (stratified, at least one episode per scenario per split), which is the
dataset behind our defensible headline.

\begin{table}[t]
\centering
\caption{Datasets. ``strat'' = stratified split guaranteeing every scenario in every split.}
\label{tab:datasets}
\small
\begin{tabularx}{\columnwidth}{@{}lXl@{}}
\toprule
\textbf{Dataset} & \textbf{Episodes / split} & \textbf{Notes} \\
\midrule
\texttt{ewat\_v3}       & 299; 209/45/45 (strat) & $\sim$21 steps; disk\_io 16.7\% NaN \\
\texttt{ewat\_v4}       & 375; 262/56/57 (temp.) & 47--51 steps; 4 scen.\ absent train \\
\texttt{ewat\_v4\_strat}& 270/60/45 (strat)      & \textbf{headline dataset} \\
\texttt{ewat\_rcaeval}  & 90; 30 fault types     & adapted RCAEval RE2-OB \\
\bottomrule
\end{tabularx}
\end{table}

\paragraph{Protocol and methodology corrections.}
Unless stated otherwise, H1/H3 on \texttt{ewat\_v3} use the optimised configuration
(average$+$cosine clustering, $d_{\text{proj}}=64$, margin~$2.0$, tuned logistic regression) and
are aggregated over ten seeds. Two methodological corrections apply throughout: (i) H1 silhouette
on val/test is computed with \emph{nearest-centroid} labels from the training centroids, not an
independent per-split clustering (which inflated scores); (ii) the precursor horizon $k^*$ is
selected on validation, never on test. Confidence intervals are $95\%$ bootstrap (BCa for AUROC).
We name the dataset and configuration for every reported number, because they differ
substantially between \texttt{ewat\_v3} and \texttt{ewat\_v4\_strat}.
